\subsection*{受理成功以后工作去了哪里}

沿用第七章的购买身份、输入版本与平台范围，先取一笔独立购买：Alice 确认 Q7、CNY 3000 分，平台账户初始为 10000 分，尚无收款。A 已保存 P1、C1、原输入与固定键 \code{pay-C1}。P1 为 \code{awaiting\_payment}，C1 为 \code{unconfirmed}，尚无 O1。本节要补上的记录，是 A 以后由谁继续这项工作。商品、付款关联和结果解释均从原 P/C 快照取得。

第七章已经把 \code{awaiting\_payment} 建立为可提交的业务阶段：A 先保住受理事实及平台映射，E 与 R 才能在退出以后继续衔接。第一部分单事务内部的 \code{pending} 随事务提交终态或回滚，现在两次提交之间的等待必须持久存在。本章为这份等待记录安排持续执行路径。

一种直接实现是在 A 以后向内存队列放入 C1。进程恰好在 A 已提交、入队尚未发生时退出，重启后的内存队列为空。P1 确实存在，任务却失去了运行机会。把入队挪到提交以前，会出现队列先交给工作者、A 随后回滚的历史，工作者甚至可能为一项未受理购买调用平台。图~\ref{fig:ch08-memory-gap} 展开了这两个退出位置。

\CEightFigure{ch08-memory-gap}{内存唤醒与受理提交分开时，两种先后顺序都留下缝隙}{fig:ch08-memory-gap}

我们在购买库中增加 \code{recovery\_tasks}，每项支付保留一条恢复任务。任务 T1 引用 C1，唯一约束使一项 C1 至多有一项对应任务。任务与 P1、C1 一起写入 A，提交后才允许工作者执行外部动作。任务身份表达哪份工作需要继续，支付身份表达平台应认作哪次付款，两者有固定关联。

\begin{BookCode}
-- 受理事务 A 内，:charge 来自支付 INSERT 的 RETURNING
INSERT INTO recovery_tasks
  (task_id, charge_id, phase, lease_epoch, next_run_at)
VALUES
  (:new_task_id, :charge, 'ready', 0, clock_timestamp())
RETURNING task_id, charge_id;
\end{BookCode}

\code{task\_id} 由受理服务分配，\code{charge\_id} 指向已保存的平台映射。\code{phase} 由调度与完成事务更新：\code{ready} 表示可以在指定时间以后领取，\code{leased} 表示某次领取留有归属记录，\code{done} 表示本章的支付恢复已取得并保存终态。\code{lease\_epoch} 是领取代次，初值 0，每次新领取递增，供后续写入识别旧工作者。\code{next\_run\_at} 是下次允许调度的数据库时刻，由受理或重新安排任务的事务写入，扫描器据此寻找候选。上述插入要求恰好返回一行，返回的任务编号及 C 才进入调度关系，示例中它们为 T1 与 C1。

受理的前半段沿用第 7.4 节：先按已有 \(S\) 的保存版本比较并恢复，新身份才进入首次业务判断。新增任务只发生在这笔首次受理事务内，实际返回的 P 传给 C，实际返回的 C 再传给 T。已有 \(S\) 对应原来的 P/C/T，恢复者读取其中保存的报价、解释版本和支付输入。工作者更换以后，继续执行这项已经受理的付款。相同 \(S\) 混用输入时，服务通过原绑定仲裁一次提交，保护已保存输入，并明确返回异参冲突。

任务插入零行、主键或 C 唯一冲突，以及其他约束异常，均回滚整个 A。失败事务中的候选 P/C/T 随之撤销，发送入口只接受已提交的关联。随后按原 \(S\) 重读，已经提交的绑定走恢复，无法对应的任务冲突进入核查。A 的提交确认丢失时，也按原 \(S\) 读取完整关联，再判断受理结局。对已提交却缺任务的旧记录，上线新增机制以前要按未决 C 完整补建并核对，核对完成后才开放发送入口。本章后续历史从三者共同提交的规则开始。

\CEightFigure{ch08-original01}{A 共同保存购买、支付映射与恢复任务}{fig:ch08-acceptance-task}

若 A 回滚，P1、C1、T1 一起消失。若 A 提交，三者一起可见。服务在提交确认后退出，扫描器可以找到 T1。内存唤醒可以缩短等待时间，扫描器则提供持续发现路径。外部平台 E 和本地保存结论 R 各自提交，新增任务保住的是它们之间继续工作的依据。

\subsection*{从一条任务记录到持续恢复}

扫描器周期性查找已经到期的 \code{ready} 任务，以及租期已经到期的 \code{leased} 任务。候选查询只提供待核对的编号，多个扫描器可以同时找到 T1，领取事务再决定谁取得本轮资格。一次查询读到 \code{ready}，随后另一个工作者已经领取，是允许发生的交错。

任务持续存在，还要保证它的状态能够再次被扫描。W1 查询平台得到处理中，随后将 T1 重新安排到稍后的数据库时刻，清除本轮持有者，令 \code{phase=ready}。若 W1 在这次安排以前退出，原领取记录到期后可被接管。若安排已经提交、确认丢失，扫描器按已保存时间继续，W1 重新读取 T1，再判断自己是否拥有原领取。

这条链构成这里的持久化工作流：业务步骤跨越多次进程运行，每次继续所需的身份、进度和下一次运行依据都保存在数据库中。\footnote{这里用持久化工作流描述跨进程保留工作依据的组织方式，单条任务配合明确的状态迁移即可实现当前恢复链。采用通用工作流引擎时，要核对它保存哪些步骤结果、怎样重放，以及外部调用使用什么固定身份。}记录中的步骤说明接下来怎样取得证据，C1 的支付输入与平台身份取自 A。工作者再次启动，就从这份保存的映射继续。

E1 后来关闭，购买入口依据第三章的会话准入规则拒绝该会话下的购买请求。已经受理的 P1/C1 由服务端继续恢复，经授权的业务事实查询也读取原映射。若客户端把同一付款改用另一会话或动作发送，完整 \(S\) 将改变，服务会把它识别为另一项购买，可能与原付款分别生效。

另一种做法扫描所有未决 C，再恢复缺失任务。条件为 \code{evidence\_state = unconfirmed}。扫描过程需要覆盖所有满足条件的记录，协调并发插入，保持已越过游标的旧未决记录可被发现，并在创建任务后重新核对支付状态。只扫描最近一天，或者扫描一次便永久推进时间游标，会遗漏更早的未决 C。本章选择 A 中共同提交任务，使覆盖依据直接落在受理事务上。

任务出错也要保留继续路径。数据库暂时不可达，工作者放弃本轮，等待原租期到期。平台暂时不可达，在持有调度资格时保存下一次运行时间。固定支付发送窗口已经结束，则把 \code{work\_mode} 从 \code{automatic} 改为 \code{check}。这个字段由恢复调度事务写入，表示当前继续方式。\code{check} 继续安排可用的查询、账单核对或有责任归属的人工核查，未决任务因此继续保存在扫描范围内。

\CEightFigure{ch08-task-cycle}{领取、重排与到期接管都保留未决工作，证据终态才结束任务}{fig:ch08-task-cycle}

\subsection*{哪项事实允许任务结束}

任务结束需要与 C1 的终态相接。取得成功证据时，R 同时保存 X1、建立唯一 O1、在 C1 保存 \code{payment\_result}、将 P1 改为 \code{paid}，并令 T1 为 \code{done}。最终拒绝证据须确认原固定身份此前没有成功，并排除已受理、在途及未来尝试成功。同一 R 在 C1 保存拒绝依据与原结果，将 P1 改为 \code{payment\_rejected} 并结束 T1，此时保持没有 O1。处理中、超时、某次调用过期，以及仅未来关闭而既往未明，都继续保留未决恢复。

在这套规则下，每个已提交且为 \code{unconfirmed} 的 C1 都有一项非 \code{done} 的 T1。T1 为 \code{done} 时，C1 已有本章定义的终态证据。先改 T1 为 \code{done}，再另开事务保存 R，会留下任务已经消失于扫描范围而 P1 等待付款的历史。先提交 R 再安排任务结束，可以保住支付事实，却留下冗余恢复。共同提交使任务生命周期直接对应本地结论。

这个覆盖关系可以逐个迁移检查。A 建立未决 C1 时同时建立 T1。领取、续期与再次安排保留 T1 的非完成状态。进程退出只留下原记录。终态合流把 C1 与 T1 一起结束。只要所有删除和状态写入都通过这些受保护的路径，受理后的工作便有可追查的持久位置。
